% !TEX program = xelatex
% ============================================================
%  Arabic Technical Report Template — Nibras Center
%  Compile with:
%    xelatex main.tex
%    xelatex main.tex
% ============================================================
%  Features:
%    - A4, 11pt, article class
%    - Title block with report number
%    - Executive summary, TOC, numbered sections
%    - Manual thebibliography (no BibTeX)
%    - Fancy headers with report title
%    - 1.5 line spacing
% ============================================================

\documentclass[11pt,a4paper]{article}

% ---- Core packages ----
\usepackage{url,amsmath}
\usepackage{bidi}
\usepackage[no-math]{fontspec}
\usepackage[quiet]{polyglossia}
\usepackage{graphicx}
\usepackage{enumitem}
\usepackage{titlesec}
\usepackage{indentfirst}
\usepackage[titles]{tocloft}
\usepackage{etoolbox}
\usepackage[labelsep=period]{caption}
\usepackage{float}
\usepackage{longtable}
\usepackage{setspace}
\usepackage{fancyhdr}
\setlength{\headheight}{14pt}

% ---- Fonts ----
\setmainfont[Scale=1]{NotoKufiArabic-Regular.ttf}
\newfontfamily\arabicfont[Script=Arabic,Scale=0.95]{NotoKufiArabic-Regular.ttf}
\newfontfamily\englishfont[Scale=0.95]{TeX Gyre Termes}
\setmonofont{NotoKufiArabic-Regular.ttf}
\let\arabicfonttt\ttfamily

% ---- Languages ----
\setdefaultlanguage[calendar=gregorian,hijricorrection=1]{arabic}
\setotherlanguage[variant=british]{english}

% ---- Arabic caption names ----
\addto\captionsarabic{%
  \renewcommand{\contentsname}{الفهرس}%
  \renewcommand{\figurename}{الشكل}%
  \renewcommand{\tablename}{الجدول}%
  \renewcommand{\refname}{المراجع}%
}

% ---- Spacing ----
\setstretch{1.5}
\setlength{\parindent}{.5in}
\setcounter{secnumdepth}{5}

% ---- Section formatting (right-aligned) ----
\titleformat{\section}
  {\normalfont\Large\bfseries}{\filright\thesection}{1ex}{\filright}
\titleformat{\subsection}
  {\normalfont\large\bfseries}{\filright\thesubsection}{1ex}{\filright}
\titleformat{\subsubsection}
  {\normalfont\normalsize\bfseries}{\filright\thesubsubsection}{1ex}{\filright}

% ---- Report metadata ----
\newcommand{\reporttitle}{عنوان التقرير الفني}
\newcommand{\reportnumber}{رقم التقرير: 001/2026}

% ---- Fancy headers ----
\pagestyle{fancy}
\fancyhf{}
\fancyhead[R]{\small\reporttitle}
\fancyhead[L]{\small\reportnumber}
\fancyfoot[C]{\thepage}
\renewcommand{\headrulewidth}{0.4pt}

% ---- Hyperref (load last) ----
\usepackage[breaklinks,hidelinks]{hyperref}

\begin{document}

% ---- Title block ----
\begin{center}
    {\LARGE\textbf{\reporttitle}}\\[1cm]
    {\large إعداد: اسم المنظمة}\\[0.5cm]
    {\large \today}\\[0.5cm]
    {\large \reportnumber}
\end{center}

\vspace{1cm}

% ---- Executive summary ----
\section*{الملخص التنفيذي}
\addcontentsline{toc}{section}{الملخص التنفيذي}

% TODO: اكتب الملخص التنفيذي هنا. يلخص الغرض والنتائج والتوصيات الرئيسية.

% ---- Table of contents ----
\tableofcontents
\newpage

% ---- Numbered sections ----

\section{المقدمة}

% TODO: اكتب مقدمة التقرير هنا. وضّح الخلفية والأهداف والنطاق.

\section{الخلفية}

% TODO: اكتب خلفية الموضوع هنا.

\section{المنهجية}

% TODO: اكتب منهجية العمل هنا.

% ---- Example table (commented out) ----
% \begin{table}[H]
% \centering
% \caption{عنوان الجدول}
% \begin{tabular}{|l|l|l|}
% \hline
% العمود 1 & العمود 2 & العمود 3 \\ \hline
% بيانات & بيانات & بيانات \\ \hline
% \end{tabular}
% \end{table}

\section{النتائج}

% TODO: اكتب النتائج هنا.

% ---- Example figure (commented out) ----
% \begin{figure}[H]
% \centering
% \includegraphics[width=0.8\textwidth]{figure.png}
% \caption{عنوان الشكل}
% \end{figure}

\section{المناقشة}

% TODO: اكتب مناقشة النتائج هنا.

\section{الاستنتاجات}

% TODO: اكتب الاستنتاجات هنا.

\section{التوصيات}

% TODO: اكتب التوصيات هنا.

% ---- References (manual thebibliography) ----
\begin{thebibliography}{99}

\bibitem{ref1}
اسم المؤلف. \textit{عنوان المرجع}. اسم الناشر، المدينة، 2024.

\bibitem{ref2}
اسم المؤلف. ``عنوان المقال.'' \textit{اسم المجلة}، المجلد 1، 2024، ص. 1--10.

\end{thebibliography}

\end{document}
